\chapter{Branch Prediction Subsystem}

\section{Overview}

One of the biggest performance bottlenecks in modern processors is the presence of branch
instructions. Every conditional branch introduces uncertainty because the processor does not
immediately know which instruction should be fetched next. Waiting until the branch is resolved
would stall the pipeline and significantly reduce performance.

To avoid these stalls, Sargantana performs \textbf{speculative execution} by predicting the outcome of
branches before they are executed. If the prediction is correct, instruction execution continues
without interruption. If the prediction is incorrect, speculative instructions are discarded and the
processor restarts from the correct target address.

Unlike a simple five-stage processor, Sargantana performs branch prediction during the
\textbf{Fetch-1 (F1)} stage so that the following stages remain continuously supplied with instructions.

\begin{rtlrefs}
  \rtlitem{rtl/datapath/rtl/if_stage_1/rtl/if_stage_1.sv}
  \rtlitem{rtl/datapath/rtl/if_stage_1/rtl/branch_predictor.sv}
\end{rtlrefs}

\begin{figure}[htbp]
  \centering
  \includegraphics[width=0.68\textwidth]{ch05_btb_bht.png}
  \caption{Program Counter fields feeding the TAG, FSM, and TARGET structures of the branch predictor}
  \label{fig:btb-bht}
\end{figure}

\section{Speculative Execution}

Speculative execution allows instructions located after a predicted branch to begin moving
through the pipeline before the branch outcome is known.

The branch predictor selects one of two possible next PCs:
\begin{itemize}
  \item Sequential PC (PC + 4)
  \item Predicted branch target
\end{itemize}

The Fetch stage immediately begins fetching instructions from the selected address.
If the prediction later proves correct, execution continues normally.
If incorrect, the processor:
\begin{itemize}
  \item flushes incorrect pipeline entries,
  \item restores architectural state,
  \item redirects the PC to the correct target.
\end{itemize}

This recovery mechanism minimizes branch penalties while maintaining correctness.

\begin{rtlrefs}
  \rtlitem{rtl/datapath/rtl/if_stage_1/rtl/if_stage_1.sv}
  \rtlitem{rtl/control_unit/rtl/control_unit.sv}
\end{rtlrefs}

\section{Branch Target Buffer (BTB)}

The Branch Target Buffer stores previously encountered branch instructions together with their
destination addresses.

Whenever a new PC enters Fetch-1, the BTB is searched.
If a matching entry exists, the BTB immediately supplies the predicted destination address
instead of waiting for the branch instruction to execute.

Therefore, BTB answers the question:

\begin{quote}
\textit{``Where should execution go?''}
\end{quote}

Its primary purpose is reducing control hazards caused by indirect and conditional branches.
Within Sargantana, BTB functionality is managed by the branch prediction subsystem inside the
Fetch-1 stage.

\begin{rtlrefs}
  \rtlitem{rtl/datapath/rtl/if_stage_1/rtl/branch_predictor.sv}
\end{rtlrefs}

\section{Branch History Table (BHT)}

Predicting only the target address is insufficient because the processor must also determine
whether the branch will actually be taken.

The Branch History Table records the recent behavior of each branch instruction.
Instead of storing target addresses, it stores prediction information indicating whether the
branch is likely to be:
\begin{itemize}
  \item Taken
  \item Not Taken
\end{itemize}

Sargantana implements this functionality through the Bimodal Predictor, which uses saturating
prediction counters that are continuously updated using previous execution history.

The BHT therefore answers:

\begin{quote}
\textit{``Will this branch be taken?''}
\end{quote}

Together, the BTB and BHT determine both the destination and the likelihood of branching.

\begin{rtlrefs}
  \rtlitem{rtl/datapath/rtl/if_stage_1/rtl/bimodal_predictor.sv}
  \rtlitem{rtl/datapath/rtl/if_stage_1/rtl/branch_predictor.sv}
\end{rtlrefs}

\section{Return Address Stack (RAS)}

Function calls and returns occur very frequently in software.
Predicting return instructions using only the BTB is often inaccurate because a single function
may be called from many different locations.

To improve prediction accuracy, Sargantana implements a Return Address Stack (RAS).

Operation:
\begin{itemize}
  \item CALL instruction $\rightarrow$ push return address.
  \item RETURN instruction $\rightarrow$ pop return address.
  \item popped address becomes predicted next PC.
\end{itemize}

This significantly improves prediction accuracy for nested function calls.

The Return Address Stack is implemented as an independent hardware structure inside the
Fetch stage.

\begin{rtlrefs}
  \rtlitem{rtl/datapath/rtl/if_stage_1/rtl/return_address_stack.sv}
\end{rtlrefs}

\section{Prediction Flow}

The overall prediction process performed during instruction fetch is summarized below.

\begin{enumerate}
  \item Current Program Counter enters Fetch-1.
  \item Branch Predictor receives the PC.
  \item BTB checks whether the PC corresponds to a known branch.
  \item BHT predicts whether the branch will be taken.
  \item RAS provides return addresses for function returns when required.
  \item Predictor selects the next PC.
  \item Fetch-2 retrieves instructions from the predicted address.
  \item Later execution stages verify the correctness of the prediction.
  \item If incorrect, Recovery Logic flushes the pipeline and redirects execution.
\end{enumerate}

This mechanism allows Sargantana to maintain a high instruction throughput while minimizing
branch penalties.

\begin{rtlrefs}
  \rtlitem{rtl/datapath/rtl/if_stage_1/rtl/if_stage_1.sv}
  \rtlitem{rtl/datapath/rtl/if_stage_1/rtl/branch_predictor.sv}
  \rtlitem{rtl/control_unit/rtl/control_unit.sv}
\end{rtlrefs}

\section{Recovery from Mispredictions}

A branch prediction is considered complete only after the branch instruction executes.
If the prediction matches the actual result, execution proceeds normally.
Otherwise, the processor performs a recovery sequence:

\begin{itemize}
  \item Flush speculative instructions.
  \item Restore the correct architectural state.
  \item Update predictor history.
  \item Redirect the Program Counter.
  \item Resume instruction fetch.
\end{itemize}

The recovery logic works together with:
\begin{itemize}
  \item Instruction Queue
  \item Rename Tables
  \item Free Lists
  \item Graduation List
  \item Control Unit
\end{itemize}

to ensure that incorrectly executed speculative instructions never modify the committed
architectural state.

\begin{rtlrefs}
  \rtlitem{rtl/control_unit/rtl/control_unit.sv}
  \rtlitem{rtl/datapath/rtl/ir_stage/rtl/instruction_queue.sv}
  \rtlitem{rtl/datapath/rtl/ir_stage/rtl/rename_table.sv}
  \rtlitem{rtl/datapath/rtl/ir_stage/rtl/free_list.sv}
  \rtlitem{rtl/datapath/rtl/wb_stage/rtl/graduation_list.sv}
\end{rtlrefs}
